\documentclass[10pt, letterpaper]{article}

% --- 基础设置 ---
\usepackage[margin=1in, footskip=0.3in]{geometry}
\usepackage[T1]{fontenc}
\usepackage[utf8]{inputenc}
\usepackage{amsmath,amssymb}
\usepackage{booktabs}
\usepackage{caption}
\usepackage{microtype}
\usepackage{subcaption}
\usepackage{graphicx}
\usepackage[table]{xcolor} % 必须加 [table] 才能用 \rowcolor

% --- 现代排版：消除缩进，增加段距 ---
\usepackage{parskip} 

% --- 字体建议：使用标准自带但更美观的字体 ---
\usepackage{mathpazo} % 使用 Palatino 字体，比默认字体更稳重
\usepackage[scaled=0.95]{inconsolata} % 极佳的代码字体

% --- 颜色定义 (Caltech/Professional Blue) ---
\definecolor{CaltechBlue}{RGB}{0, 94, 184}
\definecolor{CodeGray}{gray}{0.96}
\definecolor{CommentGreen}{RGB}{0, 128, 0}

\usepackage{hyperref}
\hypersetup{colorlinks=true, linkcolor=CaltechBlue, urlcolor=CaltechBlue, citecolor=CaltechBlue}

% --- 代码块美化 (仿 VS Code 浅色风格) ---
\usepackage{listings}
\lstset{
    basicstyle=\ttfamily\small,
    backgroundcolor=\color{CodeGray},
    keywordstyle=\color{blue}\bfseries,
    commentstyle=\color{CommentGreen}\itshape,
    stringstyle=\color{red!70!black},
    breaklines=true,
    frame=none, % 现代感：不要边框，只用背景色
    xleftmargin=10pt,
    framesep=5pt,
    showstringspaces=false,
    columns=flexible
}

% 2. 改进 Listings 设置
\lstset{
  language=C++,
  basicstyle=\ttfamily\scriptsize, % 如果 10pt 还是挤，改用 scriptsize
  keywordstyle=\color{blue!70!black}\bfseries,
  commentstyle=\color{green!50!black}\itshape,
  stringstyle=\color{red!60!black},
  breaklines=true,
  breakatwhitespace=true, % 只在空格处换行，防止截断变量名
  frame=single,
  frametrule=0.5pt,
  rulecolor=\color{gray!40},
  backgroundcolor=\color{gray!3},
  showstringspaces=false,
  columns=flexible, % 允许字符间距微调
  keepspaces=true,
  xleftmargin=1.5em, % 增加内边距
  xrightmargin=1.5em,
  aboveskip=1em,
  belowskip=1em,
}

% 3. 针对表格的建议：如果表格还是太宽
% 在 \begin{table} 下面加一句 \small 或 \footnotesize

\lstdefinestyle{python}{
  language=Python,
  basicstyle=\ttfamily\footnotesize,
  keywordstyle=\color{blue!70!black}\bfseries,
  commentstyle=\color{green!50!black}\itshape,
  stringstyle=\color{red!60!black},
  breaklines=true,
  frame=single,
  framerule=0.5pt,
  rulecolor=\color{gray!50},
  backgroundcolor=\color{gray!5},
  numbers=none,
  showstringspaces=false,
  columns=flexible,
  xleftmargin=1em,
  xrightmargin=1em,
}

% --- 标题美化 ---
\usepackage{titlesec}
\titleformat{\section}{\large\bfseries\color{CaltechBlue}}{\thesection}{1em}{}[\titlerule]

\title{\vspace{-1cm} \textbf{\texttt{torch.linalg.eigh}: Analysis of the Performance Cliff at Matrix Size $n=32$} \\ 
      \large (and cuSOLVER Integration Fix)}
\author{Sijing Du \\ \small California Institute of Technology}
\date{February 2026}


\begin{document}
\maketitle

\begin{abstract}
\texttt{torch.linalg.eigh} suffers a catastrophic performance
cliff when the matrix size crosses $n = 32 \to 33$ for batched inputs
on CUDA GPU. The root cause is PyTorch's size-dependent dispatch logic, which only calls
the newest batched cuSOLVER API (\texttt{cusolverDnXsyevBatched}) for
$n \le 32$, but returns to a sequential \textit{for-loop} for $n>32$. We demonstrate the problem, propose a fix to pyTorch's dispatching heuristics that calls cuSOLVER \texttt{cusolverDnXsyevBatched} for all matrix size $n$---matching
JAX's fix (PR~\#31375)---which eliminates the cliff entirely. The fix requires changing only \textbf{a single line} of
dispatch condition in PyTorch's C++ source code.
\end{abstract}

\paragraph{Test environment.}
All benchmarks: NVIDIA GeForce RTX~4080, CUDA~12.8,
PyTorch~2.10.0+cu128, \texttt{float64}, Ubuntu~22.04.5~LTS (WSL2).
Timing via CUDA events (3~warmup, 5~timed repetitions).


%% ================================================================
\section{Observation in default PyTorch \texttt{eigh} performance}
\label{sec:observation}
%% ================================================================

\texttt{torch.linalg.eigh} applied to a batch of $(B, n, n)$ symmetric
matrices exhibits smooth scaling for $n \le 32$, then an abrupt
slowdown at $n = 33$ (Table~\ref{tab:default}).

\begin{table}[h]
\centering
\caption{Default \texttt{torch.linalg.eigh} timing (ms) vs.\ matrix size $n$ ($B = 128$, \texttt{float64}).}
\label{tab:default}
\begin{tabular}{r rr}
\toprule
$n$ & time (ms) & ratio to $n\!=\!32$ \\
\midrule
     8  &    0.32 &         \\
    16  &    0.66 &         \\
    24  &    1.22 &         \\
    28  &    1.44 &         \\
    30  &    1.56 &         \\
    31  &    1.77 &         \\
    32  &    1.54 & 1.0$\times$   \\
\midrule
    33  &   137.7 & 89$\times$   \\
    34  &   145.8 & 95$\times$   \\
    40  &   160.9 & 104$\times$   \\
    48  &   180.6 & 117$\times$   \\
    56  &   204.4 & 133$\times$   \\
    64  &   241.4 & 157$\times$   \\
    80  &   387.9 & 252$\times$   \\
    96  &   458.0 & 297$\times$   \\
   128  &   605.1 & 393$\times$   \\
   256  &   714.2 & 464$\times$   \\
   400  &  1113.7 & 723$\times$   \\
   512  &  1382.3 & 897$\times$   \\
   513  &  1455.7 & 945$\times$   \\
   600  &  1793.8 & 1165$\times$  \\
  1024  &  3555.8 & 2309$\times$  \\
\bottomrule
\end{tabular}
\end{table}

\subsection{Root cause: PyTorch's dispatch heuristics}

The dispatch lives in
\texttt{aten/src/ATen/native/cuda/linalg/BatchLinearAlgebraLib.cpp},
function \texttt{linalg\_eigh\_cusolver} (line $\sim$1614):

\begin{lstlisting}
// Default pytorch eigh solver dispatch logic
void linalg_eigh_cusolver(const Tensor& eigenvalues, const Tensor& eigenvectors, const Tensor& infos, bool upper, bool compute_eigenvectors) {
#if defined(USE_ROCM)
  ... //(code for AMD GPUs)
#else // not USE_ROCM
  if (batchCount(eigenvectors) > 1 && eigenvectors.size(-1) <= 32) {
    // Use syevjBatched for batched matrix operation when matrix size <= 32
    // See https://github.com/pytorch/pytorch/pull/53040#issuecomment-788264724
    linalg_eigh_cusolver_syevj_batched(eigenvalues, eigenvectors, infos, upper, compute_eigenvectors);
  } else if (eigenvectors.scalar_type() == at::kFloat && eigenvectors.size(-1) >= 32 && eigenvectors.size(-1) <= 512) {
    // syevj is better than syevd for float32 dtype and matrix sizes 32x32 - 512x512
    // See https://github.com/pytorch/pytorch/pull/53040#issuecomment-788264724
    linalg_eigh_cusolver_syevj(eigenvalues, eigenvectors, infos, upper, compute_eigenvectors);
  } else {
    linalg_eigh_cusolver_syevd(eigenvalues, eigenvectors, infos, upper, compute_eigenvectors);
  }
#endif
}
\end{lstlisting}

The branching heuristics:
\begin{itemize}
    \item \textbf{$n \le 32$, batch $> 1$}: correctly routes to
    \texttt{linalg\_eigh\_cusolver\_syevj\_batched()}, which (since
    PR~\#155695, cuSOLVER~$\ge$~11701) calls
    \texttt{cusolverDnXsyevBatched}---a new batched API alternative to the legacy \texttt{cusolverDn<t>syevjBatched}. (We find that this is fast even for batched matrices of size $n>32$.)

    \item \textbf{Otherwise}: falls through to the
    \texttt{else~if} or \texttt{else} branches, which call single-matrix
    wrappers (\texttt{syevj} or \texttt{syevd}) in a \emph{loop}. Each
    matrix is solved sequentially with a separate kernel launch. This is
    where the 80--120$\times$ cliff comes from.
\end{itemize}

The $n \le 32$ gate for batched input is a leftover from the poor performance of the old \texttt{syevjBatched}
Jacobi API (see PR~\#53040), which was replaced by the newer/faster 
\texttt{cusolverDnXsyevBatched} in PR~\#155695 (June~2025). The
top-level dispatch condition was never updated.

\newpage

\section{Solution: dispatch updated cuSOLVER batched
  \texttt{eigh} for $n>32$,
  as in JAX PR~\#31375}
\label{sec:fix}
%% ================================================================
It turns out that cuSOLVER has already implemented an optimized batched \texttt{eigh} function for all matrix size $n$, which is just never called for $n>32$ in pyTorch. To get rid of the performance cliff at $n=32$, one just needs to change the cuSOLVER dispatching heuristics in pyTorch's C++ source file that interfaces with cuSOLVER API.

\subsection{The fix}

\texttt{cusolverDnXsyevBatched} (available since cuSOLVER~11.7.1 /
CUDA~12.6) is a newer batched eigensolver (divide-and-conquer, or Jacobi for $n\le32$) alternative to older \texttt{cusolverDn<t>syevjBatched}. PyTorch already wraps it and uses it
for $n \le 32$. The fix is to remove the $n \le 32$ gate:

\begin{lstlisting}
// Proposed change in linalg_eigh_cusolver:
// cusolverDnXsyevBatched works for all n when batch size > 1.
// See jax-ml/jax#31375 for precedent.
void linalg_eigh_cusolver(const Tensor& eigenvalues, const Tensor& eigenvectors, const Tensor& infos, bool upper, bool compute_eigenvectors) {
#if defined(USE_ROCM)
  ... //(code for AMD GPUs)
#else // not USE_ROCM
  if (batchCount(eigenvectors) > 1) {
    // FIX: Use syevjBatched for batched matrix operation for all matrix size n
    linalg_eigh_cusolver_syevj_batched(eigenvalues, eigenvectors, infos, upper, compute_eigenvectors);
  } else if (eigenvectors.scalar_type() == at::kFloat && eigenvectors.size(-1) >= 32 && eigenvectors.size(-1) <= 512) {
    // syevj is better than syevd for float32 dtype and matrix sizes 32x32 - 512x512
    // See https://github.com/pytorch/pytorch/pull/53040#issuecomment-788264724
    linalg_eigh_cusolver_syevj(eigenvalues, eigenvectors, infos, upper, compute_eigenvectors);
  } else {
    linalg_eigh_cusolver_syevd(eigenvalues, eigenvectors, infos, upper, compute_eigenvectors);
  }
#endif
}
\end{lstlisting}

This is to route all batched \texttt{eigh} through
\texttt{cusolverDnXsyevBatched}, replacing the per-matrix
\texttt{syevj}\&\texttt{syevd} loop for $n>32$. This fix is probably similar to the approach JAX took in
\href{https://github.com/jax-ml/jax/pull/31375}{jax-ml/jax\#31375}
(merged September~2025).

\subsection{Verification}

To verify the fix, we compare the performance of \texttt{torch.linalg.eigh} before and after applying the fix.

\paragraph{Performance.}
Tables~\ref{tab:f64_B1}--\ref{tab:f32_B256} show the results across
batch sizes $B \in \{1, 32, 64, 128, 256\}$ and dtypes
\texttt{float64}/\texttt{float32}.

%% ---------- float64 ----------



\begin{table}[h]
\centering
\caption{Default vs.\ fixed dispatch ($B = 128$, \texttt{float64}).}
\label{tab:f64_B128}
\begin{tabular}{r rrr}
\toprule
$n$ & default (ms) & fixed (ms) & speedup \\
\midrule
     8  &    0.32 &    0.31 & 1.0$\times$ \\
    16  &    0.66 &    0.63 & 1.1$\times$ \\
    24  &    1.22 &    1.19 & 1.0$\times$ \\
    28  &    1.44 &    1.49 & 1.0$\times$ \\
    30  &    1.56 &    1.66 & 0.9$\times$ \\
    31  &    1.77 &    1.64 & 1.1$\times$ \\
    32  &    1.54 &    1.59 & 1.0$\times$ \\
\midrule
    33  &   137.7 &    3.04 & \textbf{45$\times$} \\
    34  &   145.8 &    2.63 & \textbf{55$\times$} \\
    40  &   160.9 &    3.19 & \textbf{50$\times$} \\
    48  &   180.6 &    3.47 & \textbf{52$\times$} \\
    56  &   204.4 &    4.12 & \textbf{50$\times$} \\
    64  &   241.4 &    4.67 & \textbf{52$\times$} \\
    80  &   387.9 &    10.5 & \textbf{37$\times$} \\
    96  &   458.0 &    12.9 & \textbf{36$\times$} \\
   128  &   605.1 &    18.9 & \textbf{32$\times$} \\
   256  &   714.2 &    68.3 & \textbf{10$\times$} \\
   400  &  1113.7 &   198.9 & \textbf{6$\times$} \\
   512  &  1382.3 &   322.5 & \textbf{4$\times$} \\
   513  &  1455.7 &   391.0 & \textbf{4$\times$} \\
   600  &  1793.8 &   547.6 & \textbf{3$\times$} \\
  1024  &  3555.8 &  1962.0 & \textbf{1.8$\times$} \\
\bottomrule
\end{tabular}
\end{table}





\subsection{Summary}

\begin{enumerate}
    \item The $n = 32 \to 33$ performance cliff is \textbf{eliminated}---timing scales smoothly across all~$n$.
    \item \textbf{Huge speedup} for $n > 32$ (at $B=128$, \texttt{float64}), with no
    regression for $n \le 32$.
    \item The fix is a \textbf{one-line change}, using an API PyTorch already wraps.
    \item This matches what
    \href{https://github.com/jax-ml/jax/pull/31375}{JAX shipped in
    September~2025}.
\end{enumerate}


%% ================================================================
\newpage
\appendix
\section{Minimal Reproducer}
\label{app:reproducer}

\begin{lstlisting}[style=python]
import torch
torch.manual_seed(42)
device = "cuda"
B = 128
dtype = torch.float64
print(f'torch version: {torch.__version__}')
print(f'B={B}, dtype={dtype}, matrix size n')
w = 8
print(f"  {'n':>{w}} | {'eigh (ms)':>{w}} ")
print("-" * (2 * w + 9))
for n in [1,2,4,8,16,32,33,34,64,96,128]:
    x = torch.randn(B, n, n, device=device, dtype=dtype)
    a = (x + x.mT) / 2
    # warmup
    for _ in range(1):
        torch.linalg.eigh(a)
    torch.cuda.synchronize()
    
    start = torch.cuda.Event(enable_timing=True)
    end = torch.cuda.Event(enable_timing=True)
    start.record()
    for _ in range(5):
        eigvals, eigvecs = torch.linalg.eigh(a)
    end.record()
    torch.cuda.synchronize()
    t = start.elapsed_time(end) / 5
    marker = " <-- n=32 threshold" if n == 32 else ""
    print(f"  {n:{w}d} | {t:{w}.2f} ms{marker}")
\end{lstlisting}

\section{Timing Data}
Output of the minimal reproducer before the fix (torch version: 2.10.0+cu126):
\begin{lstlisting}[style=python]
torch version: 2.10.0+cu126
B=128, dtype=torch.float64, matrix size n
         n | eigh (ms) 
-------------------------
         1 |     0.12 ms
         2 |     0.33 ms
         4 |     0.20 ms
         8 |     0.30 ms
        16 |     0.62 ms
        32 |     1.65 ms <-- n=32 threshold
        33 |   137.13 ms
        34 |   139.70 ms
        64 |   228.77 ms
        96 |   457.44 ms
       128 |   580.13 ms
\end{lstlisting}
Output of the minimal reproducer after the fix (customized torch build):
\begin{lstlisting}[style=python]
torch version: 2.12.0a0+git67c428c
B=128, dtype=torch.float64, matrix size n
         n | eigh (ms) 
-------------------------
         1 |     0.13 ms
         2 |     0.35 ms
         4 |     0.23 ms
         8 |     0.29 ms
        16 |     0.62 ms
        32 |     1.71 ms <-- n=32 threshold
        33 |     2.98 ms
        34 |     2.72 ms
        64 |     5.01 ms
        96 |    14.23 ms
       128 |    19.08 ms
\end{lstlisting}
%% ---------- float64 ----------
\begin{table}[h]
\centering
\caption{Default vs.\ fixed dispatch ($B = 1$, \texttt{float64}).}
\label{tab:f64_B1}
\begin{tabular}{r rrr}
\toprule
$n$ & default (ms) & fixed (ms) & speedup \\
\midrule
     8  &    0.31 &    0.29 & 1.1$\times$ \\
    16  &    0.54 &    0.44 & 1.2$\times$ \\
    24  &    0.59 &    0.51 & 1.2$\times$ \\
    28  &    0.63 &    0.74 & 0.9$\times$ \\
    30  &    0.65 &    0.58 & 1.1$\times$ \\
    31  &    0.79 &    0.60 & 1.3$\times$ \\
    32  &    0.72 &    0.63 & 1.1$\times$ \\
\midrule
    33  &    1.31 &    1.18 & 1.1$\times$ \\
    34  &    1.28 &    1.22 & 1.1$\times$ \\
    40  &    1.43 &    1.43 & 1.0$\times$ \\
    48  &    1.88 &    1.57 & 1.2$\times$ \\
    56  &    1.95 &    1.73 & 1.1$\times$ \\
    64  &    2.00 &    1.94 & 1.0$\times$ \\
    80  &    3.48 &    3.17 & 1.1$\times$ \\
    96  &    3.73 &    3.66 & 1.0$\times$ \\
   128  &    4.87 &    4.84 & 1.0$\times$ \\
   256  &    5.78 &    5.66 & 1.0$\times$ \\
   400  &    8.98 &    8.76 & 1.0$\times$ \\
   512  &    11.3 &    11.2 & 1.0$\times$ \\
   513  &    11.8 &    11.7 & 1.0$\times$ \\
   600  &    14.3 &    14.0 & 1.0$\times$ \\
  1024  &    28.3 &    28.2 & 1.0$\times$ \\
\bottomrule
\end{tabular}
\end{table}

\begin{table}[h]
\centering
\caption{Default vs.\ fixed dispatch ($B = 32$, \texttt{float64}).}
\label{tab:f64_B32}
\begin{tabular}{r rrr}
\toprule
$n$ & default (ms) & fixed (ms) & speedup \\
\midrule
     8  &    0.26 &    0.29 & 0.9$\times$ \\
    16  &    0.38 &    0.42 & 0.9$\times$ \\
    24  &    0.59 &    0.59 & 1.0$\times$ \\
    28  &    0.62 &    0.67 & 0.9$\times$ \\
    30  &    0.62 &    0.63 & 1.0$\times$ \\
    31  &    0.67 &    0.98 & 0.7$\times$ \\
    32  &    0.65 &    0.73 & 0.9$\times$ \\
\midrule
    33  &    37.0 &    1.43 & \textbf{26$\times$} \\
    34  &    36.8 &    1.24 & \textbf{30$\times$} \\
    40  &    42.1 &    1.39 & \textbf{30$\times$} \\
    48  &    47.6 &    1.51 & \textbf{32$\times$} \\
    56  &    54.7 &    1.75 & \textbf{31$\times$} \\
    64  &    64.1 &    1.83 & \textbf{35$\times$} \\
    80  &   100.6 &    4.38 & \textbf{23$\times$} \\
    96  &   116.0 &    5.04 & \textbf{23$\times$} \\
   128  &   147.4 &    7.09 & \textbf{21$\times$} \\
   256  &   178.7 &    26.9 & \textbf{7$\times$} \\
   400  &   278.7 &    66.5 & \textbf{4$\times$} \\
   512  &   348.9 &    90.7 & \textbf{4$\times$} \\
   513  &   367.8 &   108.3 & \textbf{3$\times$} \\
   600  &   457.2 &   149.9 & \textbf{3$\times$} \\
  1024  &   892.1 &   505.5 & \textbf{1.8$\times$} \\
\bottomrule
\end{tabular}
\end{table}

\begin{table}[h]
\centering
\caption{Default vs.\ fixed dispatch ($B = 64$, \texttt{float64}).}
\label{tab:f64_B64}
\begin{tabular}{r rrr}
\toprule
$n$ & default (ms) & fixed (ms) & speedup \\
\midrule
     8  &    0.27 &    0.27 & 1.0$\times$ \\
    16  &    0.41 &    0.48 & 0.8$\times$ \\
    24  &    0.67 &    0.68 & 1.0$\times$ \\
    28  &    0.78 &    0.88 & 0.9$\times$ \\
    30  &    0.88 &    0.81 & 1.1$\times$ \\
    31  &    0.99 &    0.83 & 1.2$\times$ \\
    32  &    0.89 &    0.82 & 1.1$\times$ \\
\midrule
    33  &    71.1 &    1.66 & \textbf{43$\times$} \\
    34  &    73.7 &    1.50 & \textbf{49$\times$} \\
    40  &    80.5 &    1.63 & \textbf{49$\times$} \\
    48  &    91.0 &    1.89 & \textbf{48$\times$} \\
    56  &   105.3 &    2.18 & \textbf{48$\times$} \\
    64  &   117.5 &    2.61 & \textbf{45$\times$} \\
    80  &   194.9 &    5.75 & \textbf{34$\times$} \\
    96  &   228.0 &    6.84 & \textbf{33$\times$} \\
   128  &   295.8 &    9.85 & \textbf{30$\times$} \\
   256  &   356.2 &    44.7 & \textbf{8$\times$} \\
   400  &   555.3 &   109.2 & \textbf{5$\times$} \\
   512  &   698.0 &   159.8 & \textbf{4$\times$} \\
   513  &   737.4 &   189.5 & \textbf{4$\times$} \\
   600  &   916.1 &   275.1 & \textbf{3$\times$} \\
  1024  &  1785.2 &   977.7 & \textbf{1.8$\times$} \\
\bottomrule
\end{tabular}
\end{table}
%% ---------- float32 ----------

\begin{table}[h]
\centering
\caption{Default vs.\ fixed dispatch ($B = 1$, \texttt{float32}).}
\label{tab:f32_B1}
\begin{tabular}{r rrr}
\toprule
$n$ & default (ms) & fixed (ms) & speedup \\
\midrule
   8  & 0.26  & 0.27  & 1.0$\times$ \\
  16  & 0.27  & 0.29  & 0.9$\times$ \\
  24  & 0.32  & 0.34  & 0.9$\times$ \\
  28  & 0.32  & 0.33  & 1.0$\times$ \\
  30  & 0.31  & 0.30  & 1.0$\times$ \\
  31  & 0.31  & 0.43  & 0.7$\times$ \\
  32  & 0.75  & 0.93  & 0.8$\times$ \\
\midrule
  33  & 1.04  & 1.72  & 0.6$\times$ \\
  34  & 1.08  & 1.15  & 0.9$\times$ \\
  40  & 1.25  & 1.24  & 1.0$\times$ \\
  48  & 1.53  & 1.12  & 1.4$\times$ \\
  56  & 1.34  & 1.13  & 1.2$\times$ \\
  64  & 1.27  & 1.13  & 1.1$\times$ \\
  80  & 1.77  & 1.83  & 1.0$\times$ \\
  96  & 1.77  & 1.82  & 1.0$\times$ \\
 128  & 2.31  & 2.44  & 0.9$\times$ \\
 256  & 4.92  & 4.91  & 1.0$\times$ \\
 400  & 8.99  & 9.65  & 0.9$\times$ \\
 512  & 12.55 & 12.28 & 1.0$\times$ \\
 513  & 3.65  & 3.59  & 1.0$\times$ \\
 600  & 4.42  & 4.53  & 1.0$\times$ \\
1024  & 7.76  & 7.91  & 1.0$\times$ \\
\bottomrule
\end{tabular}
\end{table}

\begin{table}[h]
\centering
\caption{Default vs.\ fixed dispatch ($B = 32$, \texttt{float32}).}
\label{tab:f32_B32}
\begin{tabular}{r rrr}
\toprule
$n$ & default (ms) & fixed (ms) & speedup \\
\midrule
   8  & 0.18  & 0.21  & 0.9$\times$ \\
  16  & 0.24  & 0.23  & 1.0$\times$ \\
  24  & 0.30  & 0.35  & 0.9$\times$ \\
  28  & 0.28  & 0.37  & 0.8$\times$ \\
  30  & 0.30  & 0.41  & 0.7$\times$ \\
  31  & 0.32  & 0.32  & 1.0$\times$ \\
  32  & 0.28  & 0.48  & 0.6$\times$ \\
\midrule
  33  & 30.0  & 0.95  & \textbf{31$\times$} \\
  34  & 28.6  & 0.49  & \textbf{59$\times$} \\
  40  & 29.4  & 0.50  & \textbf{59$\times$} \\
  48  & 30.8  & 0.51  & \textbf{61$\times$} \\
  56  & 33.6  & 0.55  & \textbf{61$\times$} \\
  64  & 33.2  & 0.61  & \textbf{55$\times$} \\
  80  & 45.8  & 0.98  & \textbf{47$\times$} \\
  96  & 45.8  & 1.09  & \textbf{42$\times$} \\
 128  & 66.9  & 1.38  & \textbf{48$\times$} \\
 256  & 156.6 & 4.98  & \textbf{31$\times$} \\
 400  & 286.8 & 24.0  & \textbf{12$\times$} \\
 512  & 370.5 & 28.3  & \textbf{13$\times$} \\
 513  & 118.0 & 46.6  & \textbf{2.5$\times$} \\
 600  & 143.3 & 50.7  & \textbf{2.8$\times$} \\
1024  & 249.6 & 121.5 & \textbf{2.1$\times$} \\
\bottomrule
\end{tabular}
\end{table}

\begin{table}[h]
\centering
\caption{Default vs.\ fixed dispatch ($B = 64$, \texttt{float32}).}
\label{tab:f32_B64}
\begin{tabular}{r rrr}
\toprule
$n$ & default (ms) & fixed (ms) & speedup \\
\midrule
     8  &    0.20 &    0.18 & 1.1$\times$ \\
    16  &    0.27 &    0.24 & 1.1$\times$ \\
    24  &    0.29 &    0.28 & 1.1$\times$ \\
    28  &    0.31 &    0.54 & 0.6$\times$ \\
    30  &    0.31 &    0.32 & 1.0$\times$ \\
    31  &    0.30 &    0.33 & 0.9$\times$ \\
    32  &    0.30 &    0.32 & 1.0$\times$ \\
\midrule
    33  &    59.3 &    0.83 & \textbf{71$\times$} \\
    34  &    58.1 &    0.52 & \textbf{111$\times$} \\
    40  &    59.9 &    0.52 & \textbf{115$\times$} \\
    48  &    62.7 &    0.57 & \textbf{110$\times$} \\
    56  &    64.6 &    0.63 & \textbf{103$\times$} \\
    64  &    67.6 &    0.67 & \textbf{101$\times$} \\
    80  &    93.1 &    1.12 & \textbf{83$\times$} \\
    96  &   101.9 &    1.30 & \textbf{78$\times$} \\
   128  &   137.5 &    1.83 & \textbf{75$\times$} \\
   256  &   307.2 &    6.36 & \textbf{48$\times$} \\
   400  &   587.1 &    31.3 & \textbf{19$\times$} \\
   512  &   754.4 &    38.9 & \textbf{19$\times$} \\
   513  &   237.8 &    60.0 & \textbf{4$\times$} \\
   600  &   282.3 &    76.9 & \textbf{4$\times$} \\
  1024  &   523.4 &   258.4 & \textbf{2$\times$} \\
\bottomrule
\end{tabular}
\end{table}

\begin{table}[h]
\centering
\caption{Default vs.\ fixed dispatch ($B = 256$, \texttt{float32}).}
\label{tab:f32_B256}
\begin{tabular}{r rrr}
\toprule
$n$ & default (ms) & fixed (ms) & speedup \\
\midrule
     8  &    0.20 &    0.20 & 1.0$\times$ \\
    16  &    0.29 &    0.30 & 0.9$\times$ \\
    24  &    0.47 &    0.50 & 0.9$\times$ \\
    28  &    0.57 &    0.62 & 0.9$\times$ \\
    30  &    0.61 &    0.59 & 1.0$\times$ \\
    31  &    0.60 &    0.60 & 1.0$\times$ \\
    32  &    0.61 &    0.64 & 1.0$\times$ \\
\midrule
    33  &   232.6 &    1.37 & \textbf{170$\times$} \\
    34  &   255.5 &    1.02 & \textbf{249$\times$} \\
    40  &   253.5 &    1.16 & \textbf{218$\times$} \\
    48  &   272.8 &    1.32 & \textbf{206$\times$} \\
    56  &   280.4 &    1.62 & \textbf{173$\times$} \\
    64  &   329.6 &    1.85 & \textbf{178$\times$} \\
    80  &   399.0 &    3.28 & \textbf{122$\times$} \\
    96  &   399.0 &    4.26 & \textbf{94$\times$} \\
   128  &   565.0 &    6.11 & \textbf{93$\times$} \\
   256  &  1249.6 &    25.6 & \textbf{49$\times$} \\
   400  &  2472.4 &   135.2 & \textbf{18$\times$} \\
   512  &  3108.4 &   214.2 & \textbf{15$\times$} \\
   513  &   956.5 &   257.5 & \textbf{4$\times$} \\
   600  &  1130.2 &   371.8 & \textbf{3$\times$} \\
  1024  &  2045.9 &  1158.7 & \textbf{1.8$\times$} \\
\bottomrule
\end{tabular}
\end{table}

\end{document}
